\documentclass[11pt,a4paper]{scrartcl}

\usepackage{amssymb}
\usepackage{amsmath}

\usepackage[utf8]{inputenc}
\usepackage[T1]{fontenc}
\newcommand{\ats}[1]{\par\noindent\textit{Atsakymas:} #1}
\usepackage{cancel}
\usepackage{tikz}
\usepackage{lmodern} 
\usepackage{amsmath,amssymb,amsthm}
\usepackage{graphicx}
\usepackage[dvipsnames,svgnames]{xcolor}
\usepackage{hyperref}
\usepackage{mathrsfs}
\usepackage[shortlabels]{enumitem}
\usepackage[obeyFinal,textsize=scriptsize,shadow,loadshadowlibrary]{todonotes}
\usepackage{mathtools}
\usepackage{microtype}

%%% Paragraph layout
\setlength{\parindent}{0pt}
\setlength{\parskip}{0.6\baselineskip}

%%% Colored sections
\renewcommand*{\sectionformat}%
  {\color{purple}\S\thesection\enskip}
\renewcommand*{\subsectionformat}%
  {\color{purple}\S\thesubsection\enskip}
\renewcommand*{\subsubsectionformat}%
  {\color{purple}\S\thesubsubsection\enskip}
\addtokomafont{paragraph}{\color{orange!35!black}\P\ }
\KOMAoptions{numbers=noenddot}

%%% Title
\addtokomafont{subtitle}{\Large}
\setkomafont{author}{\Large\scshape}
\setkomafont{date}{\Large\normalsize}
% Neigiamas vertikalus tarpas, kuris pakelia dokumento antraštę į viršų. 
% Galite keisti -2.5cm į -3cm (aukščiau) arba -1.5cm (žemiau).
\titlehead{\vspace*{-7.5cm}} 

% PRIDĖTA: Automatiškai perrašome datą į mokyklos pavadinimą
\AtBeginDocument{%
  \date{\normalsize Vilniaus licėjus}
}

%%% Page setup
\usepackage[headsepline]{scrlayer-scrpage}
\renewcommand{\headfont}{}
\addtolength{\textheight}{3.14cm}
\setlength{\footskip}{0.5in}
\setlength{\headsep}{10pt}
% Puslapio antraštė turi rodyti tik konkrečius dokumento duomenis.
% Nustatome juos aiškiai, kad neatsirastų "\@author" / "\@title" arba senų reikšmių.
\makeatletter
\AtBeginDocument{%
  \ihead{\footnotesize\textbf{Andrius Berniukevičius}}%
  \ohead{\footnotesize\textbf{\@title}}%
}
\makeatother
\automark{section}
\chead{}
\cfoot{\pagemark}

%%% Macros
\providecommand{\ol}{\overline}
\providecommand{\ul}{\underline}
\providecommand{\wt}{\widetilde}
\providecommand{\wh}{\widehat}
\providecommand{\eps}{\varepsilon}
\providecommand{\half}{\frac{1}{2}}
\providecommand{\inv}{^{-1}}
\newcommand{\dang}{\measuredangle} %% Directed angle
\newcommand{\vocab}[1]{\textbf{\color{blue}\sffamily #1}}
\providecommand{\CC}{\mathbb C}
\providecommand{\FF}{\mathbb F}
\providecommand{\NN}{\mathbb N}
\providecommand{\QQ}{\mathbb Q}
\providecommand{\RR}{\mathbb R}
\providecommand{\ZZ}{\mathbb Z}
\providecommand{\ts}{\textsuperscript}
\providecommand{\dg}{^\circ}
\providecommand{\ii}{\item}
\providecommand{\defeq}{\coloneqq}
\DeclareMathOperator*{\lcm}{lcm}
\DeclareMathOperator*{\argmin}{arg min}
\DeclareMathOperator*{\argmax}{arg max}
\providecommand{\hrulebar}{\par
\hspace{\fill}\rule{0.95\linewidth}{.7pt}\hspace{\fill}
\par\nointerlineskip \vspace{\baselineskip}}

%%% Optional colored theorems
\usepackage{thmtools}
\usepackage[framemethod=TikZ]{mdframed}
\mdfdefinestyle{mdbluebox}{%
  roundcorner=10pt,
  linewidth=1pt,
  skipabove=12pt,
  innerbottommargin=9pt,
  skipbelow=2pt,
  linecolor=blue,
  nobreak=true,
  backgroundcolor=TealBlue!5,
}
\declaretheoremstyle[
  headfont=\sffamily\bfseries\color{MidnightBlue},
  mdframed={style=mdbluebox},
  headpunct={\\[3pt]},
  postheadspace={0pt}
]{thmbluebox}
\mdfdefinestyle{mdredbox}{%
  linewidth=0.5pt,
  skipabove=12pt,
  frametitleaboveskip=5pt,
  frametitlebelowskip=0pt,
  skipbelow=2pt,
  frametitlefont=\bfseries,
  innertopmargin=4pt,
  innerbottommargin=8pt,
  nobreak=true,
  backgroundcolor=Salmon!5,
  linecolor=RawSienna,
}
\declaretheoremstyle[
  headfont=\bfseries\color{RawSienna},
  mdframed={style=mdredbox},
  headpunct={\\[3pt]},
  postheadspace={0pt},
]{thmredbox}
\mdfdefinestyle{mdgreenbox}{%
  skipabove=8pt,
  linewidth=2pt,
  rightline=false,
  leftline=true,
  topline=false,
  bottomline=false,
  linecolor=ForestGreen,
  backgroundcolor=ForestGreen!5,
}
\declaretheoremstyle[
  headfont=\bfseries\sffamily\color{ForestGreen!70!black},
  bodyfont=\normalfont,
  spaceabove=2pt,
  spacebelow=1pt,
  mdframed={style=mdgreenbox},
  headpunct={ --- },
]{thmgreenbox}
\mdfdefinestyle{mdblackbox}{%
  skipabove=8pt,
  linewidth=3pt,
  rightline=false,
  leftline=true,
  topline=false,
  bottomline=false,
  linecolor=black,
  backgroundcolor=RedViolet!5!gray!5,
}
\declaretheoremstyle[
  headfont=\bfseries,
  bodyfont=\normalfont\small,
  spaceabove=0pt,
  spacebelow=0pt,
  mdframed={style=mdblackbox}
]{thmblackbox}

%% Aplinkos su rėmeliais (thmtools)
\declaretheorem[style=thmbluebox,name=Teorema]{theorem}
\declaretheorem[style=thmbluebox,name=Lema]{lemma}
\declaretheorem[style=thmbluebox,name=Savybė]{property}
\declaretheorem[style=thmbluebox,name=Teiginys]{proposition}
\declaretheorem[style=thmbluebox,name=Išvada]{corollary}
\declaretheorem[style=thmbluebox,name=Teorema,numbered=no]{theorem*}
\declaretheorem[style=thmbluebox,name=Lema,numbered=no]{lemma*}
\declaretheorem[style=thmbluebox,name=Teiginys,numbered=no]{proposition*}
\declaretheorem[style=thmbluebox,name=Išvada,numbered=no]{corollary*}
\declaretheorem[style=thmgreenbox,name=Algoritmas]{algorithm}
\declaretheorem[style=thmgreenbox,name=Algoritmas,numbered=no]{algorithm*}
\declaretheorem[style=thmgreenbox,name=Teiginys]{claim}
\declaretheorem[style=thmgreenbox,name=Teiginys,numbered=no]{claim*}
\declaretheorem[style=thmredbox,name=Pavyzdys]{example}
\declaretheorem[style=thmredbox,name=Pavyzdys,numbered=no]{example*}
\declaretheorem[style=thmblackbox,name=Pastaba]{remark}
\declaretheorem[style=thmblackbox,name=Pastaba,numbered=no]{remark*}

%% Standartinės amsthm aplinkos (Apibrėžimai ir užduotys)
\theoremstyle{definition}
\ifcsname conjecture\endcsname\else\newtheorem{conjecture}{Hipotezė}\fi
\ifcsname definition\endcsname\else\newtheorem{definition}{Apibrėžimas}\fi
\ifcsname fact\endcsname\else\newtheorem{fact}{Faktas}\fi
\ifcsname answer\endcsname\else\newtheorem{answer}{Atsakymas}\fi
\ifcsname ques\endcsname\else\newtheorem{ques}{Klausimas}\fi
\ifcsname exercise\endcsname\else\newtheorem{exercise}{Užduotis}\fi
\ifcsname problem\endcsname\else\newtheorem{problem}{Uždavinys}[section]\fi
\ifcsname conjecture*\endcsname\else\newtheorem*{conjecture*}{Hipotezė}\fi
\ifcsname definition*\endcsname\else\newtheorem*{definition*}{Apibrėžimas}\fi
\ifcsname fact*\endcsname\else\newtheorem*{fact*}{Faktas}\fi
\ifcsname answer*\endcsname\else\newtheorem*{answer*}{Atsakymas}\fi
\ifcsname ques*\endcsname\else\newtheorem*{ques*}{Klausimas}\fi
\ifcsname exercise*\endcsname\else\newtheorem*{exercise*}{Užduotis}\fi
\ifcsname problem*\endcsname\else\newtheorem*{problem*}{Uždavinys}\fi

\newcounter{answerssection}
\newcommand{\answerssection}{%
  \setcounter{answerssection}{\value{section}}%
  \section{Atsakymai}%
  \setcounter{problem}{0}%
  \renewcommand{\theproblem}{\arabic{answerssection}.\arabic{problem}}%
}

%% GALUTINIS SPRENDIMAS PROOF APLINKAI
%% --- PRADŽIA: Įrodymo aplinkos sutvarkymas ---
\makeatletter

% 1. Užtikriname, kad žodis būtų "Įrodymas", net jei "babel" paketas bando jį perrašyti
\AtBeginDocument{%
  \renewcommand{\proofname}{Įrodymas}%
  \@ifpackageloaded{babel}{%
    \addto\captionslithuanian{\renewcommand{\proofname}{Įrodymas}}%
  }{}%
}

% 2. Priverstinai pašaliname scrartcl klasės bold šriftą ir paliekame TIK italic
% 3. Sujungiame su tuščiu kvadratėliu dešiniajame kampe.
\renewcommand{\qedsymbol}{\ensuremath{\Box}}
\renewenvironment{proof}[1][\proofname]{\par
  \pushQED{\qed}%
  \normalfont \topsep6\p@\@plus6\p@\relax
  \trivlist
  \item[\hskip\labelsep
        \normalfont\itshape % Priverčia naudoti tik pasvirą šriftą, kaip nuotraukoje
    #1\@addpunct{.}]\ignorespaces
}{%
  \popQED\endtrivlist\@endpefalse
}
\newenvironment{solution}{\par
  \normalfont \topsep6\p@\@plus6\p@\relax
  \trivlist
  \item[\hskip\labelsep
        \normalfont\itshape
    \textit{Sprendimas}\@addpunct{.}]\ignorespaces
}{%
  \endtrivlist\@endpefalse
}
\newcommand{\ans}[1]{\par\noindent\textit{Atsakymas:} #1}
\makeatother


\title{Algebra I: Aritmetinis ir geometrinis vidurkis}
\author{Andrius Berniukevičius}

\begin{document}

\maketitle

\section{Nelygybės esmė}

Lygčių sprendimas paprastai reiškia tikslių nežinomųjų reikšmių radimą. Nelygybėse dažnai svarbesnė ne konkreti reikšmė, o visų galimų reikšmių įvertinimas: reikia nustatyti, kokių ribų reiškinys negali peržengti.

Olimpiadinėje matematikoje dažnai prašoma įrodyti, kad reiškinys yra teigiamas, ne mažesnis už tam tikrą skaičių, arba rasti jo mažiausią ar didžiausią reikšmę. Tokiose užduotyse svarbu ne spėti atsakymą, o rasti patikimą argumentą, kuris veiktų visais leistinais atvejais.

Šiame skyriuje remsimės pagrindiniu algebros faktu: bet kokio realaus skaičiaus kvadratas yra neneigiamas t.y. $x^2 \ge 0$. Iš šio paprasto fakto išauga visa olimpiadinių nelygybių teorija.


% ==========================================
% BAZINIAI GINKLAI
% ==========================================
\section{Pradiniai nelygybių įrankiai}

Nelygybių uždaviniuose svarbu atpažinti reiškinio struktūrą ir parinkti tinkamą įvertinimo būdą. Toliau aptarsime du pagrindinius įrankius: kvadratų išskyrimą ir aritmetinio bei geometrinio vidurkių nelygybę.

\subsection{Kvadratų išskyrimo savybės}
Jei bet koks skaičius pakeltas kvadratu yra $\ge 0$, vadinasi, ir bet koks skirtumas pakeltas kvadratu yra $\ge 0$:
$$ (x - y)^2 \ge 0 $$
Atskliaudę šią formulę gauname $x^2 - 2xy + y^2 \ge 0$. Perkėlę $-2xy$ į kitą pusę, gauname patį paprasčiausią, bet dažniausiai naudojamą nelygybių įrankį:
$$ {x^2 + y^2 \ge 2xy} $$
Ši taisyklė sako: dviejų skaičių kvadratų suma \textit{visada} yra didesnė arba lygi jų dvigubai sandaugai. Lygybė ($=$) galioja tik vieninteliu atveju – kai $x = y$.

\subsection{AM-GM nelygybė}
Tai yra pati garsiausia nelygybė visoje mokyklinėje matematikoje. Jos pavadinimas šifruojamas kaip \textit{Aritmetinis vidurkis $\ge$ Geometrinis vidurkis} (Arithmetic Mean $\ge$ Geometric Mean).
\begin{remark*}
AM-GM nelygybė taikoma tik teigiamiems skaičiams.
\end{remark*}

\begin{proposition}[AM-GM nelygybė dviem skaičiams]
Jei $a, b > 0$, tai
$$ \frac{a+b}{2} \geq \sqrt{ab}. $$
Lygybė pasiekiama tada ir tik tada, kai $a=b$.
\end{proposition}

\begin{proof}
Kadangi kiekvieno realiojo skaičiaus kvadratas yra neneigiamas, turime
$$ (\sqrt{a}-\sqrt{b})^2 \geq 0. $$
Atskliaudę gauname
$$ a-2\sqrt{ab}+b \geq 0, $$
todėl
$$ a+b \geq 2\sqrt{ab}. $$
Padaliję abi puses iš $2$, gauname teiginio nelygybę. Lygybė galima tik tada, kai $\sqrt{a}=\sqrt{b}$, t. y. kai $a=b$.
\end{proof}

Trijų teigiamų skaičių $a, b$ ir $c$ atveju:
$$ a + b + c \ge 3\sqrt[3]{abc} $$

Bendruoju atveju, jei $a_1, a_2, \ldots, a_n > 0$, tai
$$
\frac{a_1+a_2+\cdots+a_n}{n}
\ge
\sqrt[n]{a_1a_2\cdots a_n}.
$$
Lygybė pasiekiama tada ir tik tada, kai $a_1=a_2=\cdots=a_n$.

\textbf{Kame AM-GM yra magija?} 
Ji leidžia \textbf{sumą paversti sandauga}. Olimpiadose dažnai gauname sumą reiškinių, kurie dauginant vienas kitą suprastina. Pavyzdžiui, skaičiaus ir jam atvirkštinio skaičiaus suma:
$$ x + \frac{1}{x} \ge 2\sqrt{x \cdot \frac{1}{x}} \implies x + \frac{1}{x} \ge 2 $$

% ==========================================
% PAVYZDŽIAI
% ==========================================
\section{Sprendimų pavyzdžiai}

\begin{example}[Kvadratų atskyrimas]
Įrodykite, kad bet kokiems realiesiems skaičiams $a, b$ ir $c$ galioja nelygybė:
$$ a^2 + b^2 + c^2 \ge ab + bc + ca $$
\end{example}
\begin{proof}

Matome kvadratus ir dviejų kintamųjų sandaugas. Trūksta tik dvejetų, kad susidarytų pilni kvadratai pagal formulę $(x-y)^2$. Kadangi nelygybes galima padauginti iš teigiamo skaičiaus jų nepakeičiant, padauginkime abi nelygybės puses iš 2:
$$ 2a^2 + 2b^2 + 2c^2 \ge 2ab + 2bc + 2ca $$
Sukelkime viską į kairę pusę:
$$ 2a^2 + 2b^2 + 2c^2 - 2ab - 2bc - 2ca \ge 0 $$
Dabar „išardykime“ $2a^2$ į $a^2 + a^2$ (taip pat padarykime su $b$ ir $c$) ir pergrupupukime narius taip, kad susidarytų trys pilni kvadratai:
$$ (a^2 - 2ab + b^2) + (b^2 - 2bc + c^2) + (c^2 - 2ca + a^2) \ge 0 $$
Suskleiskime kiekvieną grupę pagal skirtumo kvadrato formulę:
$$ (a - b)^2 + (b - c)^2 + (c - a)^2 \ge 0 $$
Kadangi bet koks realaus skaičiaus kvadratas yra $\ge 0$, tai ir trijų kvadratų suma visada bus $\ge 0$. Nelygybė teisinga.
\end{proof}

\begin{example}[Atvirkštinių dydžių minimumas]
Raskite mažiausią įmanomą reiškinio $x + \frac{16}{x}$ reikšmę, kai žinoma, kad $x > 0$. Su kokia $x$ reikšme šis minimumas pasiekiamas?
\end{example}
\begin{solution}

Matome dviejų teigiamų skaičių sumą. Vieno iš jų kintamasis yra skaitiklyje, kito – vardiklyje. Tai tobulas scenarijus AM-GM nelygybei, nes šiuos skaičius sudauginus, $x$ susiprastins ir liks tik skaičius. Pritaikome AM-GM nelygybę $a + b \ge 2\sqrt{ab}$, kur $a = x$, o $b = \frac{16}{x}$:
$$ x + \frac{16}{x} \ge 2\sqrt{x \cdot \frac{16}{x}} $$
Susiprastiname sandaugą po šaknimi:
$$ x + \frac{16}{x} \ge 2\sqrt{16} $$
$$ x + \frac{16}{x} \ge 2 \cdot 4 $$
$$ x + \frac{16}{x} \ge 8 $$
Įrodėme, kad šis reiškinys \textbf{niekada nebus mažesnis už 8}. Tai ir yra jo mažiausia galima reikšmė.
AM-GM nelygybė tampa lygybe (pasiekia patį dugną, t.y. 8) tik tuo vieninteliu atveju, kai abu sudedamieji nariai yra lygūs:
$$ x = \frac{16}{x} $$
Dauginame abi puses iš $x$:
$$ x^2 = 16 $$
Kadangi sąlygoje nurodyta $x > 0$, gauname $x = 4$. Minimumas pasiekiamas, kai $x = 4$.\\
\ats{$8$, kai $x = 4$}
\end{solution}

\begin{example}[Skliaustų daugyba]
Teigiami skaičiai $a, b$ ir $c$ tenkina sąlygą. Įrodykite, kad:
$$ (a+b)(b+c)(c+a) \ge 8abc $$
\end{example}
\begin{proof}
Neskubėkite aklai atskliaudinėti šio reiškinio – gausite 8 narius ir tiesiog pasimesite. Kur kas elegantiškiau yra pritaikyti AM-GM nelygybę kiekvienam skliaustui atskirai.
Pritaikome AM-GM nelygybę sumai $(a+b)$:
$$ a+b \ge 2\sqrt{ab} $$
Lygiai taip pat pritaikome AM-GM antram ir trečiam skliaustui:
$$ b+c \ge 2\sqrt{bc} $$
$$ c+a \ge 2\sqrt{ca} $$
Kadangi visi skaičiai $a, b, c$ yra teigiami, mes galime šias tris nelygybes saugiai sudauginti tarpusavyje (kairę pusę su kaire, dešinę su dešine):
$$ (a+b)(b+c)(c+a) \ge 2\sqrt{ab} \cdot 2\sqrt{bc} \cdot 2\sqrt{ca} $$
Sutvarkome dešinę pusę sudaugindami skaičius ir sujungdami viską po viena šaknimi:
$$ (a+b)(b+c)(c+a) \ge 8\sqrt{a^2b^2c^2} $$
Ištraukus šaknį iš kvadratų, gauname tai, ką ir reikėjo įrodyti:
$$ (a+b)(b+c)(c+a) \ge 8abc $$
\end{proof}

\begin{example}[Vardiklių žudymas triuku „Pridėk ir taikyk AM-GM“]
Įrodykite, kad bet kokiems teigiamiems skaičiams $a, b, c$ galioja:
$$ \frac{a^2}{b} + \frac{b^2}{c} + \frac{c^2}{a} \ge a + b + c $$
\end{example}
\begin{proof}
Šis uždavinys iš pirmo žvilgsnio atrodo sunkiai įkandamas, nes kairėje turime trupmenas su kvadratais skaitikliuose, o dešinėje – paprastus skaičius. Jei visą kairę pusę kišime po AM-GM šaknimi, gausime $\sqrt[3]{\frac{a^2b^2c^2}{bca}} = \sqrt[3]{abc}$, kas visai nepanašu į $a+b+c$. Mums reikia „išlaisvinti“ kvadratus.
Paimkime pirmąją trupmeną $\frac{a^2}{b}$. Kaip joje atsikratyti vardiklio $b$ ir sumažinti laipsnį nuo kvadrato iki paprasto $a$? Ogi pridėjus prie jos patį skaičių $b$ ir pritaikius AM-GM šiai dviejų narių sumai!
$$ \frac{a^2}{b} + b \ge 2\sqrt{\frac{a^2}{b} \cdot b} $$
Šaknies viduje $b$ susiprastina, ir lieka $2\sqrt{a^2} = 2a$.
Gavome tobulą asmeninę nelygybę pirmajai trupmenai:
$$ \frac{a^2}{b} + b \ge 2a $$
Analogiškai pritaikome šį „gelbėjimosi“ triuką kitoms dviem trupmenoms:
$$ \frac{b^2}{c} + c \ge 2b $$
$$ \frac{c^2}{a} + a \ge 2c $$
Surenkame visas tris nelygybes ir jas sudedame stulpeliu:
$$ \left(\frac{a^2}{b} + \frac{b^2}{c} + \frac{c^2}{a}\right) + (a + b + c) \ge 2a + 2b + 2c $$
Iš abiejų nelygybės pusių atimame po vieną $(a+b+c)$. Kairėje pusėje ši suma pradingsta, o dešinėje iš dvigubos sumos lieka tik vienguba:
$$ \frac{a^2}{b} + \frac{b^2}{c} + \frac{c^2}{a} \ge a + b + c $$
Įrodyta elegantiškai ir be varginančio bendravardiklinimo.
\end{proof}

\vspace{0.5cm}

% ==========================================
% UŽDAVINIAI
% ==========================================
\newpage
\section{Uždaviniai savarankiškam sprendimui (12 iššūkių)}

\begin{enumerate}
    \renewcommand{\labelenumi}{\textbf{\theenumi.}}

    \item Įrodykite, kad bet kokiems realiesiems skaičiams $x$ ir $y$ teisinga nelygybė:
    $$ x^2 + y^2 \ge 2xy $$

    \item Raskite mažiausią įmanomą reiškinio reikšmę, kai žinoma, kad $x > 0$:
    $$ x + \frac{9}{x} $$

    \item Raskite mažiausią įmanomą reiškinio reikšmę, kai žinoma, kad $a > 0$:
    $$ 4a + \frac{25}{a} $$

    \item Įrodykite, kad bet kokiems teigiamiems skaičiams $a$ ir $b$ teisinga nelygybė:
    $$ (a+1)(b+1) \ge 4\sqrt{ab} $$
\end{enumerate}

\begin{enumerate}
    \setcounter{enumi}{4}

    \item Raskite mažiausią įmanomą reiškinio reikšmę, kai $x > 0$:
    $$ x + \frac{x}{4} + \frac{16}{x^2} $$

    \item Įrodykite, kad bet kokiems realiesiems skaičiams $a, b$ teisinga nelygybė:
    $$ a^2 + b^2 + 2 \ge 2(a + b) $$

    \item Įrodykite, kad bet kokiems teigiamiems $a, b, c, d$ teisinga:
    $$ (a+b)(c+d) \ge 4\sqrt{abcd} $$

    \item Įrodykite, kad bet kokiems teigiamiems $x, y, z$ teisinga:
    $$ \frac{x}{y} + \frac{y}{z} + \frac{z}{x} \ge 3 $$
\end{enumerate}

\begin{enumerate}
    \setcounter{enumi}{8}

    \item Įrodykite, kad teigiamiems $a, b, c$ teisinga:
    $$ \frac{a+b}{c} + \frac{b+c}{a} + \frac{c+a}{b} \ge 6 $$

    \item Žinoma, kad teigiamų skaičių suma $x + y = 1$. Įrodykite, kad:
    $$ \left(x + \frac{1}{x}\right)^2 + \left(y + \frac{1}{y}\right)^2 \ge \frac{25}{2} $$

    \item Įrodykite, kad bet kokiems teigiamiems $a, b, c$ teisinga nelygybė:
    $$ a^3 + b^3 + c^3 \ge a^2b + b^2c + c^2a $$

    \item \textbf{(Nesbito nelygybė)} Įrodykite legendinę, visame pasaulyje garsią olimpiadinę nelygybę bet kokiems teigiamiems $a, b, c$:
    $$ \frac{a}{b+c} + \frac{b}{a+c} + \frac{c}{a+b} \ge \frac{3}{2} $$
\end{enumerate}


% ==========================================
% ATSAKYMAI
% ==========================================


\end{document}